% ======================================================================
%  preamble.tex -- shared by part1.tex, part2.tex, part3.tex
%  A* Search Algorithm -- CSE200, CSE BUET
%  ENGINE: XeLaTeX only (fontspec + OpenType)
% ======================================================================

\usetheme{default}
\useinnertheme{rounded}
\setbeamertemplate{navigation symbols}{}
\setbeamercovered{invisible}

% ----------------------------------------------------------------------
%  1.  FONTS
%      text  -> beamer default  (Latin Modern Sans, beamer's own face)
%      diag  -> Latin Modern Roman  (the LaTeX default face)
%      alg   -> Latin Modern Roman  (the LaTeX default face)
%      math  -> beamer default, untouched
%
%      >>> IBM PLEX SANS: to bring it back as the text font, uncomment
%      >>> the two \set...font lines below.  (TeX Live `plex' fonts,
%      >>> found via kpathsea -- nothing has to live in fonts/.)
% ----------------------------------------------------------------------
\usepackage{fontspec}
% \setmainfont{IBMPlexSans}[Extension=.otf,
%   UprightFont=*-Regular, BoldFont=*-Bold,
%   ItalicFont=*-Italic,   BoldItalicFont=*-BoldItalic]
% \setsansfont{IBMPlexSans}[Extension=.otf,
%   UprightFont=*-Regular, BoldFont=*-Bold,
%   ItalicFont=*-Italic,   BoldItalicFont=*-BoldItalic]
\newfontfamily\diagfont{lmroman10-}[Extension=.otf,
  UprightFont=*regular, BoldFont=*bold,
  ItalicFont=*italic,   BoldItalicFont=*bolditalic]
\let\algfont\diagfont

\usepackage{amsmath,amssymb,mathtools}
\usepackage{booktabs}
\usepackage{pifont}                % \ding symbols (part 2)
\usepackage[table]{xcolor}
\usepackage{graphicx}
\usepackage{tikz}
\usetikzlibrary{arrows.meta,positioning,decorations.pathreplacing,calc,
                shapes.geometric,backgrounds,fit,shadows.blur, tikzmark}

% ----------------------------------------------------------------------
%  2.  PALETTE
% ----------------------------------------------------------------------
\definecolor{rblue}{HTML}{2E2E9C}
\definecolor{rbluedk}{HTML}{1C1C6B}
\definecolor{rbluemd}{HTML}{5A5AC4}
\definecolor{rbluelt}{HTML}{E7E7F7}
\definecolor{cbOrange}{HTML}{D97706}
\definecolor{cbOrangeLt}{HTML}{FDF0DC}
\definecolor{cbGreen}{HTML}{0B7A3B}
\definecolor{cbGreenLt}{HTML}{E4F2E9}
\definecolor{cbRed}{HTML}{B01B2E}
\definecolor{cbRedLt}{HTML}{FAE7E9}
\definecolor{crison}{HTML}{96383c}
\definecolor{neutralg}{gray}{0.46}
\definecolor{lnnum}{gray}{0.58}
\definecolor{hlyellow}{HTML}{FFE9A8}
\definecolor{gridgray}{gray}{0.82}
\definecolor{closedc}{gray}{0.80}

% grid tiles -- flat, low-saturation rose ramp: untouched -> frontier ->
% expanded, plus a charcoal wall.  S and G are solid accent tiles.
\definecolor{gfree}{HTML}{F4EDEE}
\definecolor{gopen}{HTML}{DEBBC2}
\definecolor{gclose}{HTML}{BE7F89}
\definecolor{gwall}{HTML}{3F3F46}
\definecolor{gstart}{HTML}{0B7A3B}
\definecolor{ggoal}{HTML}{B01B2E}

% soft, low-saturation tints used only by the beamer blocks -- kept apart
% from the tikz palette above so diagram colours stay untouched
\definecolor{blkInk}{HTML}{20242E}
\definecolor{blkBlueTt}{HTML}{E4E4F3}
\definecolor{blkBlueBd}{HTML}{F3F3FA}
\definecolor{blkBlueFg}{HTML}{2E2E78}
\definecolor{blkRedTt}{HTML}{F4E4E6}
\definecolor{blkRedBd}{HTML}{FBF4F4}
\definecolor{blkRedFg}{HTML}{8C2C3A}
\definecolor{blkGrnTt}{HTML}{E2EEE7}
\definecolor{blkGrnBd}{HTML}{F2F8F4}
\definecolor{blkGrnFg}{HTML}{2E6A4A}

\newcommand{\pos}[1]{\textcolor{cbGreen}{#1}}
\newcommand{\con}[1]{\textcolor{cbOrange}{#1}}
\newcommand{\bad}[1]{\textcolor{cbRed}{#1}}
\newcommand{\HL}[1]{\textcolor{rblue}{#1}}
\newcommand{\mut}[1]{\textcolor{neutralg}{#1}}

% ----------------------------------------------------------------------
%  3.  COLOUR ASSIGNMENTS
% ----------------------------------------------------------------------
\setbeamercolor{structure}{fg=rblue}
\setbeamercolor{frametitle}{bg=rblue,fg=white}
\setbeamercolor{title}{bg=rblue,fg=white}
\setbeamercolor{block title}{bg=blkBlueTt,fg=blkBlueFg}
\setbeamercolor{block body}{bg=blkBlueBd,fg=blkInk}
\setbeamercolor{block title alerted}{bg=blkRedTt,fg=blkRedFg}
\setbeamercolor{block body alerted}{bg=blkRedBd,fg=blkInk}
\setbeamercolor{block title example}{bg=blkGrnTt,fg=blkGrnFg}
\setbeamercolor{block body example}{bg=blkGrnBd,fg=blkInk}
\setbeamercolor{itemize item}{fg=rblue}
\setbeamercolor{itemize subitem}{fg=rbluemd}
\setbeamercolor{alerted text}{fg=cbRed}

\setbeamertemplate{blocks}[default]

% Square corners WITH a shadow.  Beamer offers shadow= only on the
% [rounded] variant, so the shadow is added by hand here: each block is
% captured into \blkbox, then dropped into a tikz node that carries the
% shadow.  Geometry stays beamer's [default] block, i.e. square.
%
%      >>> NO SHADOW: comment out the six \setbeamertemplate lines at the
%      >>> end of this block; [default] above then takes over again.
\usetikzlibrary{shadows}
\newsavebox{\blkbox}
\tikzset{blkshadow/.style={drop shadow={shadow xshift=1.3pt,
                                        shadow yshift=-1.3pt,
                                        opacity=0.20, fill=black!70}}}

\newcommand{\blkopen}[2]{%   #1 = title beamercolor, #2 = body beamercolor
  \par\vskip\medskipamount%
  \begin{lrbox}{\blkbox}%
  \begin{minipage}{\textwidth}%
    \begin{beamercolorbox}[colsep*=.75ex]{#1}%
      \usebeamerfont*{block title}\insertblocktitle%
    \end{beamercolorbox}%
    {\parskip0pt\par}%
    \nointerlineskip%
    \usebeamerfont{block body}%
    \begin{beamercolorbox}[colsep*=.75ex,vmode]{#2}%
      \vskip-.75ex\vbox{}%
}
\newcommand{\blkclose}{%
    \end{beamercolorbox}%
  \end{minipage}%
  \end{lrbox}%
  % The picture is 1.5ex wider than \textwidth (see inner xsep below), so
  % it is parked in a \textwidth hbox with \hss to absorb the overhang --
  % otherwise every block reports an Overfull \hbox.  The -.75ex shift
  % puts the bleed on both sides, exactly where stock beamer puts it.
  \noindent\hbox to\textwidth{\hspace*{-.75ex}%
  \begin{tikzpicture}
    % inner xsep matches the .75ex that colsep* bleeds outside the box,
    % so the node path (and therefore the shadow) covers the painted block
    \node[inner xsep=.75ex, inner ysep=0pt, outer sep=0pt, blkshadow]
          {\usebox{\blkbox}};
  \end{tikzpicture}\hss}%
  \vskip\smallskipamount%
}

\setbeamertemplate{block begin}{\blkopen{block title}{block body}}
\setbeamertemplate{block end}{\blkclose}
\setbeamertemplate{block alerted begin}{\blkopen{block title alerted}{block body alerted}}
\setbeamertemplate{block alerted end}{\blkclose}
\setbeamertemplate{block example begin}{\blkopen{block title example}{block body example}}
\setbeamertemplate{block example end}{\blkclose}

% for triangle and square
% \setbeamertemplate{itemize item}{\raisebox{0.1ex}{\scriptsize$\blacktriangleright$}}
% \setbeamertemplate{itemize subitem}{\raisebox{0.1ex}{\tiny$\blacksquare$}}
% for circle and square
\setbeamertemplate{itemize item}{\raisebox{0.1ex}{\scriptsize$\bullet$}}
\setbeamertemplate{itemize subitem}{\raisebox{0.1ex}{\tiny$\blacksquare$}}

\setbeamerfont{frametitle}{size=\large,series=\bfseries}
\setbeamerfont{block title}{size=\small,series=\bfseries}
\setbeamerfont{footline}{size=\scriptsize}
\setbeamerfont{footnote}{size=\tiny}
% beamer's default footnote template indents by \parindent 1em and puts the
% marker in a 1.8em box; both are dropped here so the note sits flush left
\setbeamertemplate{footnote}{%
  \parindent0em\noindent\raggedright
  \usebeamerfont{footnote}%
  \hbox to 0.9em{\insertfootnotemark\hfil}\insertfootnotetext\par}

% ----------------------------------------------------------------------
%  4.  FRAME TITLE
% ----------------------------------------------------------------------
\setbeamertemplate{frametitle}{%
  \nointerlineskip
  \begin{beamercolorbox}[wd=\paperwidth,ht=3.4ex,dp=1.5ex,%
                         leftskip=1.6ex,rightskip=1.6ex]{frametitle}
    \usebeamerfont{frametitle}\insertframetitle
  \end{beamercolorbox}%
}

% ----------------------------------------------------------------------
%  4b. CORNER MOTIF  (faint search lattice, right end of the title bar)
%      A 4x3 grid graph with the staircase route traced from a filled
%      start node to a double-ringed goal -- the picture A* is about.
%      Tints are mixed against rblue so it reads as a watermark on the
%      bar.  It lives INSIDE the title bar because the longest frame
%      title stops at 12.6cm of 16cm, so that strip is always free --
%      anywhere in the body area it would clash with wide text lines.
%
%      >>> TO TURN IT OFF: comment out the \addtobeamertemplate line at
%      >>> the bottom of this block (put a % in front of it). Leave
%      >>> \cornermotif defined; nothing else refers to it.
% ----------------------------------------------------------------------
\newcommand{\motifstep}{0.30cm}      % lattice pitch
\newcommand{\cornermotif}{%
  \begin{tikzpicture}[remember picture, overlay,
      mn/.style    = {circle, draw=white!38!rblue, line width=0.4pt,
                      fill=white!20!rblue, minimum size=1.5mm, inner sep=0pt},
      mon/.style   = {mn, draw=white!62!rblue, fill=white!44!rblue},
      me/.style    = {draw=white!30!rblue, line width=0.4pt},
      mpath/.style = {draw=white!68!rblue, line width=0.8pt,
                      line cap=round, line join=round}]
    \begin{scope}[shift={($(current page.north east)+(-1.45cm,-0.88cm)$)},
                  x=\motifstep, y=\motifstep]
      \foreach \c in {0,1,2,3}{\draw[me] (\c,0) -- (\c,2);}
      \foreach \r in {0,1,2}  {\draw[me] (0,\r) -- (3,\r);}
      \draw[mpath] (0,0) -- (1,0) -- (1,1) -- (2,1) -- (2,2) -- (3,2);
      \foreach \c in {0,1,2,3}{\foreach \r in {0,1,2}{\node[mn] at (\c,\r) {};}}
      \foreach \q in {(1,0),(1,1),(2,1),(2,2)}{\node[mon] at \q {};}
      \node[mn, fill=white!88!rblue, draw=white!88!rblue] at (0,0) {};
      \node[mn, draw=white!88!rblue, double, double distance=0.35pt] at (3,2) {};
    \end{scope}
  \end{tikzpicture}%
}
\addtobeamertemplate{frametitle}{}{\cornermotif}

% ----------------------------------------------------------------------
%  GLOBAL VERTICAL NUDGE
%  One knob for the whole deck: space added just under the title bar, so
%  every titled frame's content sits lower.  The frame body is centred, so
%  the visible shift is HALF of this value.  The [plain] title slide has no
%  frametitle and is untouched.  Set to 0pt to disable.
% ----------------------------------------------------------------------
\newlength{\framenudge}
\setlength{\framenudge}{4pt}
\addtobeamertemplate{frametitle}{}{\vskip\framenudge}


% ----------------------------------------------------------------------
%  5.  FOOTER
%        ----------------------------------------------------------
%        CSE200 - A* Search Algorithm   Zarif Ishmam      n / N
%      ============                                    <- progress
% ----------------------------------------------------------------------
\newcommand{\footcourse}{CSE200 - A* Search Algorithm}
\newcommand{\footauthor}{Zarif Ishmam}
\newcommand{\pbheight}{1.5pt}

\newcommand{\progressbarline}{%
  \ifnum\inserttotalframenumber>0\relax
    % capped at 1: on a first pass the total is not known yet, and in a
    % long deck n/total would make the bar wider than TeX can measure
    \pgfmathsetmacro{\pbfrac}{min(1,\insertframenumber/\inserttotalframenumber)}%
  \else
    \def\pbfrac{0}%
  \fi
  \noindent
  \textcolor{rblue}{\rule{\pbfrac\paperwidth}{\pbheight}}%
  \textcolor{rblue!14}{\rule{\dimexpr\paperwidth-\pbfrac\paperwidth\relax}{\pbheight}}%
}

\setbeamercolor{footbar}{bg=white,fg=rbluedk}
\newcommand{\footmargin}{1.8ex}
\setbeamertemplate{footline}{%
  \noindent\hspace*{\footmargin}%
  \textcolor{gray!45}{%
    \rule{\dimexpr\paperwidth-\footmargin-\footmargin\relax}{0.4pt}}%
  \hspace*{\footmargin}%
  \vskip1.5pt
  \hbox{%
    \begin{beamercolorbox}[wd=\paperwidth,ht=2.1ex,dp=0.8ex,%
                           leftskip=\footmargin,rightskip=\footmargin]{footbar}%
      \usebeamerfont{footline}\diagfont
      % zero-width boxes: the two \hfill split the strip evenly, so the
      % author sits on the exact centre line regardless of the side widths
      \makebox[0pt][l]{\textsc{\footcourse}}%
      \hfill
      \makebox[0pt][c]{\textsc{\footauthor}}%
      \hfill
      \makebox[0pt][r]{\textsc{\insertframenumber\,/\,\inserttotalframenumber}}%
    \end{beamercolorbox}%
  }%
  \vskip1pt%
  \progressbarline
  \vskip0pt%
}

% ----------------------------------------------------------------------
%  6.  CLRS-STYLE ALGORITHM  (default LaTeX face, line numbered, overlay aware)
% ----------------------------------------------------------------------
\newcounter{alglin}
\newcommand{\kw}[1]{\textcolor{rblue}{\textbf{#1}}}
\newcommand{\proc}[1]{\textsc{#1}}
\newcommand{\com}[1]{\textcolor{neutralg}{\itshape\,\(\triangleright\) #1}}
\newcommand{\idt}[1]{\hspace*{#1\dimexpr1.25em\relax}}
\newcommand{\hlon}[2]{\alt<#1>{\setlength{\fboxsep}{1.4pt}\colorbox{hlyellow}{#2}}{#2}}

% Stacked hboxes, not a tabular: LaTeX puts \@arstrut in every tabular row,
% so a hidden line still claimed a full line of height and the block never
% shrank.  As plain hboxes a hidden \aline contributes nothing at all, which
% lets the frame centre on the lines actually showing.
\newenvironment{clrs}[1][\small]
  {\begingroup\algfont#1\setcounter{alglin}{0}%
   \setlength{\parindent}{0pt}\setlength{\parskip}{0pt}%
   \begin{minipage}[t]{\linewidth}}
  {\end{minipage}\endgroup}

\newcommand{\aline}[2][1-]{%
  \stepcounter{alglin}%
  \only<#1>{\hbox{\hbox to 1.15em{\hfil\textcolor{lnnum}{\scriptsize\thealglin}}%
                  \hspace{1.3ex}\strut#2}}}

% ----------------------------------------------------------------------
%  7.  TIKZ STYLES
% ----------------------------------------------------------------------
\tikzset{
  every picture/.append style={font=\diagfont},
  vtx/.style      = {circle, draw=rbluedk, line width=0.7pt, fill=white,
                     minimum size=6.5mm, inner sep=0pt, font=\small},
  vstart/.style   = {vtx, fill=cbGreenLt, draw=cbGreen, line width=1.1pt,
                     font=\small\bfseries},
  vgoal/.style    = {vtx, fill=cbRedLt, draw=cbRed, line width=1.1pt,
                     font=\small\bfseries, double, double distance=0.6pt},
  vopen/.style    = {vtx, fill=rbluelt, draw=rblue, line width=1.1pt},
  vclosed/.style  = {vtx, fill=closedc, draw=neutralg, line width=0.9pt},
  vcur/.style     = {vtx, fill=hlyellow, draw=cbOrange, line width=1.3pt},
  edg/.style      = {draw=neutralg, line width=0.8pt},
  edgon/.style    = {draw=rblue, line width=1.6pt},
  edgpath/.style  = {draw=cbGreen, line width=2.2pt},
  edggrey/.style  = {draw=gridgray, line width=1.6pt},
  edgdead/.style  = {draw=gridgray, line width=0.6pt},
  edgorange/.style= {draw=cbOrange, line width=2.2pt},
  arr/.style      = {-{Stealth[length=2.4mm]}, line width=1.0pt},
  wlab/.style     = {font=\scriptsize, fill=white, inner sep=1pt, text=rbluedk},
  hlab/.style     = {font=\scriptsize, text=cbOrange, inner sep=1pt},
  box/.style      = {draw=rbluedk, rounded corners=2pt, fill=rbluelt,
                     minimum width=2.3cm, minimum height=0.8cm,
                     font=\small, align=center},
  boxg/.style     = {box, fill=cbGreenLt, draw=cbGreen},
  boxr/.style     = {box, fill=cbRedLt,   draw=cbRed},
  boxo/.style     = {box, fill=cbOrangeLt, draw=cbOrange},
}

% ----------------------------------------------------------------------
%  8.  THE SHARED 16 x 10 GRID
%      cell (c,r) spans [c,c+1] x [r,r+1];  centre (c+.5, r+.5)
%      S = (1,4)   G = (14,6)
%      Walls form a U that opens WEST, straight at the start, so a greedy
%      search walks into it.  Verified numbers on this map:
%        Dijkstra 137 cells / cost 19   Greedy 48 / cost 35   A* 81 / cost 19
% ----------------------------------------------------------------------
\newcommand{\gridnodefont}{\tiny\bfseries}
\newcommand{\gridnodeS}{S}
\newcommand{\gridnodeG}{G}

% --- flat-tile look: every cell is its own square, separated by a gutter.
%     No grid rules, no outer frame -- the gutters do all the separating.
\newcommand{\gpad}{0.075}               % gutter half-width, in cell units
\newcommand{\gcell}[3]{%                % #1 = colour, #2 = col, #3 = row
  \fill[#1] ({#2+\gpad},{#3+\gpad}) rectangle ({#2+1-\gpad},{#3+1-\gpad});}
\newcommand{\gletter}[4]{%              % #1 = colour, #2 = col, #3 = row, #4 = glyph
  \gcell{#1}{#2}{#3}%
  \node[text=white, font=\gridnodefont, inner sep=0pt]
        at ({#2+0.5},{#3+0.5}) {#4};}

% the wall set: a U opening WEST
\newcommand{\gwalls}{%
  \foreach \c in {4,...,12} {\gcell{gwall}{\c}{3}\gcell{gwall}{\c}{7}}%
  \foreach \r in {4,5,6} {\gcell{gwall}{12}{\r}}}

% laid down FIRST, under everything the search paints
\newcommand{\gridback}{%
  \foreach \c in {0,...,15} {\foreach \r in {0,...,9} {\gcell{gfree}{\c}{\r}}}}

% laid down LAST, over everything the search paints
\newcommand{\drawgridbase}{%
  \gwalls
  \gletter{gstart}{1}{4}{\gridnodeS}%
  \gletter{ggoal}{14}{6}{\gridnodeG}%
}

% full-size grid picture; #1 = cell size, #2 = generated body file
\newcommand{\gridpic}[2]{%
  \begin{tikzpicture}[x=#1, y=#1]
    \gridback
    \input{#2}
  \end{tikzpicture}}

% reduced grid picture for side-by-side panels
\newcommand{\gridpicsmall}[2]{%
  \begingroup
    \renewcommand{\gridnodefont}{\fontsize{3.6}{4}\selectfont\bfseries}%
    \begin{tikzpicture}[x=#1, y=#1]
      \gridback
      \input{#2}
    \end{tikzpicture}%
  \endgroup}

% tiny panels: tiles only, no letters
\newcommand{\gridpictiny}[2]{%
  \begingroup
    \renewcommand{\gpad}{0.10}%
    \renewcommand{\gridnodeS}{}\renewcommand{\gridnodeG}{}%
    \begin{tikzpicture}[x=#1, y=#1]
      \gridback
      \input{#2}
    \end{tikzpicture}%
  \endgroup}

% legend strip used under the grid frames
\newcommand{\gridlegend}{%
  \begin{tikzpicture}[x=0.34cm,y=0.34cm, baseline=0pt]
    \renewcommand{\gpad}{0.09}%
    \gcell{gclose}{0}{0}
      \node[anchor=west, font=\tiny] at (1.4,0.5) {closed};
    \gcell{gopen}{5.4}{0}
      \node[anchor=west, font=\tiny] at (6.8,0.5) {open (frontier)};
    \gcell{cbGreen}{15.0}{0}
      \node[anchor=west, font=\tiny] at (16.4,0.5) {final path};
    \gcell{gwall}{21.4}{0}
      \node[anchor=west, font=\tiny] at (22.8,0.5) {wall};
  \end{tikzpicture}}

% grey until overlay #1, bold blue on it, dark from then on
\newcommand{\lit}[2]{%
  \temporal<#1>{\textcolor{gray!50}{#2}}%
               {\textcolor{rblue}{\textbf{#2}}}%
               {\textcolor{rbluedk}{#2}}}

\newcommand{\takeaway}[1]{%
  \vfill
  \begin{beamercolorbox}[wd=\textwidth,sep=4pt,rounded=true]{block body}
    \footnotesize\textbf{Takeaway.} #1
  \end{beamercolorbox}}
